\problemname{Typo}

Gunilla wants to visit her favorite website. Sometimes, however, she mixes up certain letters, and accidentally visits a different website instead. Interestingly, she never makes the exact same mistake twice in the same address.

Given the address of her favorite website and a description of which letters she tends to mix up, output the length of the list of addresses that hackers should register to trick her into visiting their website.

\section*{Input}

The first line contains a string: the address of Gunilla's favorite website. Only the letters a-z and period appear, and the length is between 2 and 20 characters. It is not necessarily a real address or even a correctly structured address. The second line contains an integer $N$ ($1 \le N \le 13$), the number of groups of letters that Gunilla mixes up. Each of the following $N$ lines contains a string with between 2 and 26 letters, all different and chosen from a-z. The letters in each string form a group where all letters can be mixed up with each other. For example, if the string is ``boi'' it means that ``b'' can be swapped for ``o'' or ``i'', that ``o'' can be swapped for ``b'' or ``i'', and that ``i'' can be swapped for ``b'' or ``o''. No letter appears in more than one group.

\section*{Output}

An integer: the number of addresses Gunilla can end up at by making at least one mistake, but without making the exact same typo (swapping a certain letter for a certain other letter) more than once.

\section*{Scoring}
Your solution will be tested on two test groups. To earn points for a group,
you must pass all test cases in that group.

\begin{tabular}{| l | l | l |}
\hline
Group & Points & Constraints \\ \hline
1     & 30         & The answer is less than 200 and all $N$ mix-up strings have exactly 2 letters \\
2     & 30         & The answer is less than 10,000 \\ 
3     & 40         & The answer is less than 2 billion \\ \hline
\end{tabular}


\section*{Explanation of Samples}

The potential hacker addresses in the first three examples are:
\begin{verbatim}
po.pe     kartis.com      abbe
so.pe     katris.com      abbi
so.se     kartix.com      ebba
          katrix.com      ebbi
          kattix.com      ibba
          kartiz.com      ibbe
          katriz.com
          kattiz.com
\end{verbatim}
Note the absence of e.g. ``karris.com'' and ``ebbe'' where the same mistake was made twice.
